\documentclass[11pt,a4paper]{article}

% Packages
\usepackage[utf8]{inputenc}
\usepackage[T1]{fontenc}
\usepackage{amsmath,amssymb,amsfonts}
\usepackage[margin=1in]{geometry}
\usepackage{enumitem}
\usepackage{hyperref}
\usepackage{fancyhdr}
\usepackage{titlesec}
\usepackage{tocloft}

% Page style
\pagestyle{fancy}
\fancyhf{}
\fancyhead[L]{\small Lecture Notes: Permutations and Combinations}
\fancyhead[R]{\small \thepage}
\renewcommand{\headrulewidth}{0.4pt}

% Section formatting
\titleformat{\section}{\Large\bfseries}{Section \thesection}{1em}{}
\titleformat{\subsection}{\large\bfseries}{\thesubsection}{1em}{}

% Adjust table of contents appearance
\renewcommand{\cftsecfont}{\bfseries}
\renewcommand{\cftsubsecfont}{\normalfont}

\begin{document}

% Title page
\begin{titlepage}
    \centering
    \vspace*{4cm}
    
    {\Huge\bfseries Lecture Notes:\\ Permutations and Combinations\par}
    \vspace{1cm}
    {\Large Discrete Mathematics\par}
    \vspace{2cm}
    {\large \textbf{Author:} Bakdaulet Sotsial\par}
    \vspace{1cm}
    
    \vfill
    {\small Astana IT University \par}
    \vspace*{2cm}
\end{titlepage}

% Table of contents
\tableofcontents
\newpage

% Section 1: Permutations
\section{Permutations}

\subsection{Definition of Permutations}

\textbf{Definition (Permutation).} A \textit{permutation} is an ordered arrangement of objects. Given a set of distinct objects, a permutation of the set is a sequence (or ordering) of its elements. The order of the elements matters.

\textbf{Remark.} Permutations are fundamentally about \textit{arranging} objects where the order is significant. For example, the arrangement "ABC" is different from "CBA", even though they contain the same elements.

\subsection{Permutations Without Repetition}

\textbf{Theorem (Permutations of $n$ Distinct Objects).} The number of permutations of $n$ distinct objects, taken all at a time, is denoted $P(n, n)$ or $n!$, and is given by:
\begin{equation*}
    P(n, n) = n! = n \times (n-1) \times (n-2) \times \cdots \times 2 \times 1.
\end{equation*}

\textbf{Proof.} There are $n$ choices for the first position, $n-1$ remaining choices for the second position, $n-2$ for the third, and so on, down to $1$ choice for the last position. By the Product Rule, the total number of permutations is the product of these numbers.

\textbf{Theorem (Permutations of $n$ Objects Taken $r$ at a Time).} The number of permutations of $n$ distinct objects taken $r$ at a time (i.e., ordered arrangements of $r$ elements chosen from $n$), denoted $P(n, r)$ or ${}_nP_r$, is:
\begin{equation*}
    P(n, r) = \frac{n!}{(n-r)!} = n \times (n-1) \times (n-2) \times \cdots \times (n-r+1).
\end{equation*}

\textbf{Remark.} This formula follows directly from the Product Rule: $n$ choices for the first position, $n-1$ for the second, ..., $n-r+1$ for the $r$-th position.

\begin{example}
How many ways are there to arrange the letters A, B, C, D, E in a row?

\textbf{Solution.} We are arranging all $5$ distinct letters, so the number of permutations is $5! = 120$.
\end{example}

\begin{example}
In how many ways can a president, vice-president, and secretary be chosen from a club of $20$ members, if no person can hold more than one office?

\textbf{Solution.} This is a permutation of $20$ objects taken $3$ at a time, because the order (which office each person holds) matters.
\begin{equation*}
    P(20, 3) = 20 \times 19 \times 18 = 6840.
\end{equation*}
\end{example}

\begin{example}
How many different $4$-letter "words" (sequences of letters) can be formed from the letters A, B, C, D, E, F, G, H, I, J if no letter can be repeated?

\textbf{Solution.} We choose $4$ letters from $10$ and arrange them in order:
\begin{equation*}
    P(10, 4) = 10 \times 9 \times 8 \times 7 = 5040.
\end{equation*}
\end{example}

\subsection{Permutations with Repetition}

\textbf{Theorem (Permutations with Repeated Elements).} If we have $n$ objects, where there are $n_1$ indistinguishable objects of type 1, $n_2$ indistinguishable objects of type 2, $\dots$, $n_k$ indistinguishable objects of type $k$, with $n_1 + n_2 + \cdots + n_k = n$, then the number of distinct permutations of these $n$ objects is:
\begin{equation*}
    \frac{n!}{n_1! \, n_2! \, \cdots \, n_k!}.
\end{equation*}

\textbf{Intuition.} This is the multinomial coefficient. If all objects were distinct, there would be $n!$ permutations. For each type, we divide by $n_i!$ to correct for overcounting, since permuting identical items among themselves produces indistinguishable arrangements.

\begin{example}
How many distinct arrangements can be made using all the letters of the word "MISSISSIPPI"?

\textbf{Solution.} The word MISSISSIPPI has $11$ letters:
\begin{itemize}[label=--]
    \item M: 1
    \item I: 4
    \item S: 4
    \item P: 2
\end{itemize}
Total distinct permutations:
\begin{equation*}
    \frac{11!}{1! \, 4! \, 4! \, 2!} = \frac{39916800}{1 \times 24 \times 24 \times 2} = \frac{39916800}{1152} = 34650.
\end{equation*}
\end{example}

\begin{example}
How many distinct ways can we arrange $3$ red balls, $2$ blue balls, and $5$ green balls in a row?

\textbf{Solution.} Total balls $n = 3 + 2 + 5 = 10$. Number of arrangements:
\begin{equation*}
    \frac{10!}{3! \, 2! \, 5!} = \frac{3628800}{6 \times 2 \times 120} = \frac{3628800}{1440} = 2520.
\end{equation*}
\end{example}

\begin{example}
How many different signals can be made by hoisting $4$ red flags, $3$ blue flags, and $2$ yellow flags on a vertical pole, using all the flags?

\textbf{Solution.} Total flags: $9$. Distinct signals:
\begin{equation*}
    \frac{9!}{4! \, 3! \, 2!} = \frac{362880}{24 \times 6 \times 2} = \frac{362880}{288} = 1260.
\end{equation*}
\end{example}

% Section 2: Combinations
\section{Combinations}

\subsection{Definition of Combinations}

\textbf{Definition (Combination).} A \textit{combination} is an unordered selection of objects. Given a set of $n$ distinct objects, a combination of $r$ of these objects is a subset of size $r$. The order of selection does \textbf{not} matter.

\textbf{Definition (Binomial Coefficient).} The number of combinations of $n$ distinct objects taken $r$ at a time is denoted $C(n, r)$, $\binom{n}{r}$, or ${}_nC_r$, and is given by:
\begin{equation*}
    \binom{n}{r} = \frac{n!}{r! \, (n-r)!}, \quad \text{for } 0 \leq r \leq n.
\end{equation*}
By convention, $\binom{n}{0} = 1$ and $\binom{n}{n} = 1$.

\subsection{Difference Between Permutations and Combinations}

The key distinction is:
\begin{itemize}[label=$\bullet$]
    \item \textbf{Permutations:} ORDER matters. (A, B, C) and (C, B, A) are different arrangements.
    \item \textbf{Combinations:} ORDER does not matter. $\{A, B, C\}$ and $\{C, B, A\}$ represent the same subset.
\end{itemize}

\textbf{Relationship:} $\displaystyle\binom{n}{r} = \frac{P(n, r)}{r!}$. This is because each combination of $r$ objects can be arranged in $r!$ different orders to produce $r!$ distinct permutations. Thus, permutations = combinations $\times$ (ways to order the chosen elements).

\begin{example}
How many ways are there to choose a committee of $3$ people from a group of $10$?

\textbf{Solution.} Order does not matter; a committee is just a subset. Number of ways:
\begin{equation*}
    \binom{10}{3} = \frac{10!}{3! \, 7!} = \frac{10 \times 9 \times 8}{3 \times 2 \times 1} = 120.
\end{equation*}
\end{example}

\begin{example}
From a standard deck of $52$ cards, how many $5$-card hands are possible?

\textbf{Solution.} A hand is an unordered selection of $5$ cards:
\begin{equation*}
    \binom{52}{5} = \frac{52!}{5! \, 47!} = \frac{52 \times 51 \times 50 \times 49 \times 48}{5 \times 4 \times 3 \times 2 \times 1} = 2{,}598{,}960.
\end{equation*}
\end{example}

\begin{example}
A student must answer $8$ out of $10$ questions on an exam. In how many ways can the student choose the $8$ questions?

\textbf{Solution.} Order of answering does not matter; only the selection counts:
\begin{equation*}
    \binom{10}{8} = \binom{10}{2} = \frac{10 \times 9}{2 \times 1} = 45.
\end{equation*}
\end{example}

\textbf{Remark (Symmetry Property).} $\displaystyle\binom{n}{r} = \binom{n}{n-r}$. Choosing $r$ objects to include is equivalent to choosing $n-r$ objects to exclude.

\begin{example}
In a lottery, a player selects $6$ numbers from $1$ to $49$. How many possible tickets are there?

\textbf{Solution.}
\begin{equation*}
    \binom{49}{6} = \frac{49!}{6! \, 43!} = 13{,}983{,}816.
\end{equation*}
\end{example}

% Section 3: Combinations with Repetition
\section{Combinations with Repetition}

\subsection{Definition and Formula}

\textbf{Definition (Combinations with Repetition).} A \textit{combination with repetition} (or \textit{multiset combination}) is a selection of $r$ objects from a set of $n$ distinct types, where repetition of types is allowed, and the order of selection does not matter.

\textbf{Theorem (Stars and Bars).} The number of ways to choose $r$ objects from $n$ distinct types, with repetition allowed, is:
\begin{equation*}
    \binom{r + n - 1}{r} = \binom{r + n - 1}{n - 1}.
\end{equation*}

\subsection{Stars and Bars Method}

The "stars and bars" method provides an elegant proof. Represent the $r$ chosen objects as $r$ stars ($\star$). To separate the objects into $n$ categories (types), we use $n-1$ bars ($|$). The number of stars between consecutive bars indicates how many objects of each type are selected. 

The total number of symbols is $r$ stars plus $n-1$ bars = $r + n - 1$ positions. We simply choose the positions for the $r$ stars (or equivalently, the $n-1$ bars). Hence the count is $\binom{r+n-1}{r}$.

\textbf{Intuition:} This counts the number of nonnegative integer solutions to the equation $x_1 + x_2 + \cdots + x_n = r$, where $x_i$ represents the number of objects of type $i$ chosen.

\begin{example}
How many ways are there to choose $5$ scoops of ice cream from $4$ flavors (vanilla, chocolate, strawberry, mint), if multiple scoops of the same flavor are allowed and order does not matter?

\textbf{Solution.} Here $n = 4$ (types), $r = 5$ (scoops). Number of combinations with repetition:
\begin{equation*}
    \binom{5 + 4 - 1}{5} = \binom{8}{5} = \binom{8}{3} = \frac{8 \times 7 \times 6}{3 \times 2 \times 1} = 56.
\end{equation*}
\end{example}

\begin{example}
How many nonnegative integer solutions are there to the equation $x_1 + x_2 + x_3 + x_4 = 10$?

\textbf{Solution.} This is equivalent to distributing $10$ identical items into $4$ distinct boxes. Using stars and bars with $r = 10$ and $n = 4$:
\begin{equation*}
    \binom{10 + 4 - 1}{10} = \binom{13}{10} = \binom{13}{3} = \frac{13 \times 12 \times 11}{3 \times 2 \times 1} = 286.
\end{equation*}
\end{example}

\begin{example}
A bakery sells $3$ types of cookies: chocolate chip, oatmeal raisin, and sugar. In how many ways can you buy a dozen ($12$) cookies?

\textbf{Solution.} $n = 3$, $r = 12$.
\begin{equation*}
    \binom{12 + 3 - 1}{12} = \binom{14}{12} = \binom{14}{2} = \frac{14 \times 13}{2 \times 1} = 91.
\end{equation*}
\end{example}

\begin{example}
How many ways can we distribute $8$ identical balls into $5$ distinct boxes?

\textbf{Solution.} $n = 5$ (boxes), $r = 8$ (balls).
\begin{equation*}
    \binom{8 + 5 - 1}{8} = \binom{12}{8} = \binom{12}{4} = \frac{12 \times 11 \times 10 \times 9}{4 \times 3 \times 2 \times 1} = 495.
\end{equation*}
\end{example}

\textbf{Remark.} If we require that each box gets at least one ball (positive integer solutions), we first place one ball in each of the $5$ boxes, reducing the problem to distributing the remaining $r - n$ balls freely. The count becomes $\binom{(r-n) + n - 1}{r-n} = \binom{r-1}{n-1}$.

% Section 4: Binomial Theorem
\section{Binomial Theorem}

\subsection{Statement of the Binomial Theorem}

\textbf{Theorem (Binomial Theorem).} For any real numbers (or variables) $a$ and $b$, and any nonnegative integer $n$,
\begin{equation*}
    (a + b)^n = \sum_{k=0}^{n} \binom{n}{k} a^{n-k} b^{k}.
\end{equation*}

Explicitly, the expansion is:
\begin{align*}
    (a + b)^n = \binom{n}{0} a^n b^0 + \binom{n}{1} a^{n-1} b^1 + \binom{n}{2} a^{n-2} b^2 + \cdots + \binom{n}{n} a^0 b^n.
\end{align*}

The coefficients $\binom{n}{k}$ are called \textit{binomial coefficients}.

\textbf{Remark.} The Binomial Theorem generalizes familiar expansions like $(a+b)^2 = a^2 + 2ab + b^2$ and $(a+b)^3 = a^3 + 3a^2b + 3ab^2 + b^3$. It provides a systematic way to expand any power of a binomial.

\subsection{Binomial Coefficients and Pascal's Triangle}

The binomial coefficients form a triangular array known as \textbf{Pascal's Triangle}, where each entry is the sum of the two entries above it:

\[
\begin{array}{ccccccccccccc}
    & & & & & 1 & & & & & \\
    & & & & 1 & & 1 & & & & \\
    & & & 1 & & 2 & & 1 & & & \\
    & & 1 & & 3 & & 3 & & 1 & & \\
    & 1 & & 4 & & 6 & & 4 & & 1 & \\
    1 & & 5 & & 10 & & 10 & & 5 & & 1 \\
\end{array}
\]

\textbf{Theorem (Pascal's Identity).} For integers $n$ and $k$ with $1 \leq k \leq n-1$:
\begin{equation*}
    \binom{n}{k} = \binom{n-1}{k-1} + \binom{n-1}{k}.
\end{equation*}

This identity is the foundation for constructing Pascal's Triangle and can be proved algebraically or combinatorially.

\subsection{Examples of Binomial Expansions}

\begin{example}
Expand $(x + y)^5$.

\textbf{Solution.} Using the Binomial Theorem with $n = 5$:
\begin{align*}
    (x + y)^5 &= \binom{5}{0}x^5 + \binom{5}{1}x^4 y + \binom{5}{2}x^3 y^2 + \binom{5}{3}x^2 y^3 + \binom{5}{4}x y^4 + \binom{5}{5}y^5 \\
    &= x^5 + 5x^4y + 10x^3y^2 + 10x^2y^3 + 5xy^4 + y^5.
\end{align*}
\end{example}

\begin{example}
Expand $(2x - 3)^4$.

\textbf{Solution.} Let $a = 2x$ and $b = -3$.
\begin{align*}
    (2x - 3)^4 &= \sum_{k=0}^{4} \binom{4}{k} (2x)^{4-k} (-3)^k \\
    &= \binom{4}{0}(2x)^4 + \binom{4}{1}(2x)^3(-3) + \binom{4}{2}(2x)^2(-3)^2 \\
    &\quad + \binom{4}{3}(2x)(-3)^3 + \binom{4}{4}(-3)^4 \\
    &= 1 \cdot 16x^4 + 4 \cdot 8x^3 \cdot (-3) + 6 \cdot 4x^2 \cdot 9 + 4 \cdot 2x \cdot (-27) + 1 \cdot 81 \\
    &= 16x^4 - 96x^3 + 216x^2 - 216x + 81.
\end{align*}
\end{example}

\begin{example}
Find the coefficient of $x^6 y^4$ in the expansion of $(x + 2y)^{10}$.

\textbf{Solution.} The term involving $x^{10-k} y^k$ is $\binom{10}{k} x^{10-k} (2y)^k = \binom{10}{k} 2^k x^{10-k} y^k$. We need $10-k = 6$, so $k = 4$. The term is:
\begin{equation*}
    \binom{10}{4} x^{6} (2y)^4 = \binom{10}{4} \cdot 2^4 \cdot x^6 y^4 = 210 \cdot 16 \cdot x^6 y^4 = 3360 \, x^6 y^4.
\end{equation*}
Thus, the coefficient is $3360$.
\end{example}

\begin{example}
Prove the identity: $\displaystyle\sum_{k=0}^{n} \binom{n}{k} = 2^n$.

\textbf{Solution.} Set $a = 1$ and $b = 1$ in the Binomial Theorem:
\begin{equation*}
    (1 + 1)^n = \sum_{k=0}^{n} \binom{n}{k} 1^{n-k} 1^k = \sum_{k=0}^{n} \binom{n}{k}.
\end{equation*}
Since $(1+1)^n = 2^n$, the identity follows. This shows that the sum of all binomial coefficients for a given $n$ equals $2^n$, which corresponds to the total number of subsets of a set with $n$ elements.
\end{example}

\begin{example}
Find the sum $\displaystyle\sum_{k=0}^{n} (-1)^k \binom{n}{k}$.

\textbf{Solution.} Set $a = 1$ and $b = -1$ in the Binomial Theorem:
\begin{equation*}
    (1 - 1)^n = \sum_{k=0}^{n} \binom{n}{k} 1^{n-k} (-1)^k = \sum_{k=0}^{n} (-1)^k \binom{n}{k}.
\end{equation*}
Since $(1-1)^n = 0^n = 0$ for $n \geq 1$, the sum is $0$ for all $n \geq 1$. For $n = 0$, the sum is $1$.
\end{example}

\textbf{Remark.} The Binomial Theorem is a cornerstone of algebra and combinatorics. The binomial coefficients $\binom{n}{k}$ have countless appearances in mathematics, from probability theory (binomial distributions) to calculus (Taylor series) to number theory (properties of binomial coefficients modulo primes). They also connect to many combinatorial identities and provide a bridge between algebra and counting.

\end{document}